By setting $X=H^k(\Omega)$, $Y=H^{k-1}(\Omega)$ and $Z=L^1(\Omega)$, we can derive the following consequence.

\begin{Corollary}
\label{cor:ehrling-sobolev-equivalent-norms}
Assume that $\Omega$ is bounded and\[H^k(\Omega)\overset{c}{\hookrightarrow}H^{k-1}(\Omega).\]Then the following norms on $H^k(\Omega)$ are equivalent:
\begin{equation}\label{eq:sobolev-ehrling-equivalent-norms}
\begin{aligned}
\|u\|_{k,2}^2&=\sum_{|\alpha|\leq k}\|D^\alpha u\|_2^2,\\
\|u\|_{k,2,*}^2&=\sum_{|\alpha|=k}\|D^\alpha u\|_2^2+\|u\|_1^2.
\end{aligned}
\end{equation}

\end{Corollary}

We leave the proof as an exercise (Problem~\ref{ex:ehrling-sobolev-equivalent-norms}).

In the following result, we show that for the space $H_0^k$, we can leave out the term $\|u\|_1^2$ in the norms above. Moreover, we do not need to assume that $\Omega$ is bounded; it suffices that it be bounded in one direction. This result is known as Poincaré\textquotesingle s inequality.

\begin{Theorem}{Poincaré\textquotesingle s inequality}
\label{thm:poincare-zero-boundary-highest-derivatives}
Let $\Omega$ be contained in the strip $|x_1|\leq d<\infty$. Then there is a constant $c$, depending only on $k$ and $d$, such that
\begin{equation}\label{eq:poincare-zero-boundary-highest-derivative-control}
\|u\|_{k,2}^2\leq c\sum_{|\alpha|=k}\|D^\alpha u\|_2^2
\end{equation}
for every $u\in H_0^k(\Omega)$.
\end{Theorem}

\begingroup
\postdisplaypenalty=10000
\renewcommand{\endProof}{\par\nobreak\hfill$\blacksquare$\par\vspace{0.1cm}}
\begin{Proof}
We give the proof for $k=1$, the general case follows by induction. By density, it suffices to consider $u\in\mathcal{D}(\Omega)$. An integration by parts yields
\begin{equation}\label{eq:poincare-one-direction-integration-by-parts}
\begin{aligned}
\|u\|_2^2&=\int_{\Omega}1\cdot|u(\bm{x})|^2\,\mathrm{d}\bm{x}=-\int_{\Omega}x_1\dfrac{\partial}{\partial x_1}|u|^2\,\mathrm{d}\bm{x}\\
&\leq2d\|u\|_2\left\|\dfrac{\partial u}{\partial x_1}\right\|_2.
\end{aligned}
\end{equation}

\end{Proof}
\endgroup

\begin{Remark}
\label{rem:poincare-general-p}
For $p\in(1,\infty)$, an analogous result holds. In other words, if $\Omega$ satisfies the hypotheses of Theorem~\ref{thm:poincare-zero-boundary-highest-derivatives}, then there is a constant $C$ depending only on $d$, $k$ and $p$ such that
\begin{equation}\label{eq:poincare-zero-boundary-general-p}
\|u\|_{k,p}^p\leq C\sum_{|\alpha|=k}\|D^\alpha u\|_p^p
\end{equation}
for every $u\in W_0^{k,p}(\Omega)$.
\end{Remark}
